\documentclass{article}
\usepackage{amsmath, amssymb}
\begin{document}
\title{Digest: Why is Quantum Physics Based on Complex Numbers?}
\author{Felix M. Lev (Digest)}
\maketitle

\section{Introduction}
This document formally digests Felix M. Lev's paper on Galois Field Quantum Theory (GFQT). The theory substitutes the standard complex Hilbert space with a linear projective space over a finite field $GF(p^n)$. 

\section{Modular Representations and $GF(p^2)$}
The paper proves that a finite field analog of the de Sitter algebra $so(1,4)$ requires a quadratic extension $GF(p^2)$ where $p \equiv 3 \pmod 4$ to ensure that spinless irreducible representations are fully decomposable. 
This quadratic extension acts identically to the complex numbers $\mathbb{C}$, supplying the required imaginary root $i^2 = -1$.

\section{so(1,4) Algebra}
The symmetry algebra generators $M^{ab}$ satisfy:
\begin{equation}
[M^{ab}, M^{cd}] = -2i(\eta^{ac}M^{bd} + \eta^{bd}M^{ac} - \eta^{ad}M^{bc} - \eta^{bc}M^{ad})
\end{equation}
where $\eta = \text{diag}(1, -1, -1, -1, -1)$.

\section{Modular su(2) and Operators}
The algebra includes the $su(2) \times su(2)$ operators $\mathbf{J}', \mathbf{J}''$ and $R_{jk}$. The $M_{04}$ operator, analogous to the energy operator, is analyzed in the $(\mathbf{J}^2, J_3)$ basis, yielding an equidistant imaginary spectrum. This confirms that $M_{04}$ is fully decomposable without introducing any inconsistencies in $GF(p^2)$.

\end{document}
